\section{Problem 2: Two-Bar Plane Truss Design}

\subsection{Problem description}

The Two-Bar Truss problem is composed by two diagonal bars (of aluminium here) connected on a node and are fixed at two upper point (Figure 1). On the central node we applied a vertical strength. This force represents the external load that the structure must support. We look for the vertical displacement of the node.
 
\begin{figure}[h!]
    \centering
    %\includegraphics[width=0.5\textwidth]{truss_figure.png} % Placeholder
    \caption{Two-bar plane truss} 
\end{figure}

There are two design variables which are the position of the joints 1 and 2. They are defined by the two equations: $f_1$ the structural weight truss, and $f_2$ the displacement of joint 3.

\begin{equation}
f_1(x)=2\rho hx_2\sqrt{1+x_1^2}
\end{equation}
\begin{equation}
f_2(x)=\frac{Ph(1+x_1^2)^{1.5}\sqrt{1+x_1^4}}{2\sqrt{2}Ex_1^2x_2}
\end{equation}

These functions are constrained by two stress constraints functions $g_1$ (tension) and $g_2$ (compression).
\begin{equation}
g_1(x)=\frac{P(1+x_1)\sqrt{1+x_1^2}}{2\sqrt{2}x_1x_2}-\sigma_0\le0
\end{equation}
\begin{equation}
g_2(x)=\frac{P(-x_1+1)\sqrt{1+x_1^2}}{2\sqrt{2}x_1x_2}-\sigma_0\le0
\end{equation}

The variables $x_1$ and $x_2$ are calculated from the following formulas:
\begin{equation}
x_1=\frac{x}{h} \quad \text{and} \quad x_2=\frac{A}{A_{min}}
\end{equation}
\begin{equation}
0.1\le x_1\le2.25 \quad \text{and} \quad 0.5\le x_2\le2.5
\end{equation}

If the $x_1$ variable is connected to the physical sense of the truss geometry, the position boundaries are:
$x_{min}=0.1 \times h=0.254m$ and $x_{max}=2.25 \times h=5.72m$

The parameters are converted in the International System units.
\begin{itemize}
    \item $\rho=0.283lb/in^3=7833kg/m^3$
    \item $h=100 in = 2.54 m$
    \item $P=10^4lb=44482.2N$
    \item $E=3 \times 10^7lb/in^2=2.07 \times 10^{11}Pa$
    \item $\sigma_0=2 \times 10^4lb/in^2=1.38 \times 10^8Pa$
    \item $A_{min}=1in^2=6.4516 \times 10^{-4}m^2$
\end{itemize}

\subsection{Methodology / Algorithm}

\subsubsection{Penalty function for the PSO simulation}

Pymoo library already contains constraint handling with constraint violation measures during the selection process. For the multi-objective problem, it uses the Pareto ranking. To study how the penalty strength affects convergence and feasibility, explicit penalty terms will be added.

\begin{equation}
F_{penalised}=f(x)+k \times \sum(\max(0,g_i(x))^2)
\end{equation}
With: $k$ the penalty factor

Considering this equation, a penalty will be applied when the maximum between 0 and $g_i$ is $g_i$ so the constraint is violated. This is realised with the following loop where the penalty factor is fixed to 10 billion. If the penalty factor is large the algorithm avoids infeasible solutions quickly. While, if the penalty factor is small the algorithm may explore infeasible regions longer.

\begin{lstlisting}[language=Python]
for i in range(n_solutions):
    try:
        objectives, constraints = self.truss.evaluate(x[i, :])
        F[i, :] = objectives
        G[i, :] = constraints
    except ValueError:
        # Penalty for invalid solutions
        F[i, :] = [1e10, 1e10]
        G[i, :] = [1e10, 1e10]
\end{lstlisting}

\subsection{Results}

\subsubsection{PSO simulation}

The extracted results from the multi-objective Particle Swarm Optimisation are summarised below. The algorithm identified a Pareto front offering a trade-off between Weight and Displacement.

\begin{table}[h!]
\centering
\caption{Key Solutions from PSO Simulation}
\begin{tabular}{|l|c|c|c|c|}
\hline
\textbf{Solution Type} & \textbf{x1} & \textbf{x2} & \textbf{Weight (kg)} & \textbf{Displacement (m)} \\ \hline
Min Weight & 0.824 & 0.500 & 16.636 & 0.00843 \\ \hline
Min Displacement & 1.878 & 2.500 & 136.534 & 0.00108 \\ \hline
Compromise & 1.256 & 1.167 & 48.098 & 0.00292 \\ \hline
\end{tabular}
\end{table}

\begin{figure}[h!]
    \centering
    \includegraphics[width=0.8\textwidth]{pareto_front.png}
    \caption{Pareto Front: Weight vs Displacement}
\end{figure}

\subsection{Discussion}
The results show a clear trade-off. Minimising weight leads to a significant increase in displacement, while minimising displacement requires a much heavier structure. The compromise solution offers a balanced design.

\subsection{Source code}

The problem definition and Pymoo wrapper can be found in `two\_bar\_plane\_truss\_problem.py`. The simulation execution is handled by `PSO\_simulation\_multi\_obj.py`.
